\documentclass[11pt,a4paper]{article}
\usepackage[italian]{babel}
\usepackage[margin=2.5cm]{geometry}
\usepackage{parskip}
\usepackage{amsmath, amssymb, amsthm}
\usepackage{booktabs}
\usepackage{array}
\usepackage{xcolor}
\usepackage{needspace}
\usepackage{enumitem}
\usepackage{listings}
\usepackage{tikz}
\usepackage[most]{tcolorbox}
\usepackage[hidelinks]{hyperref}
\usetikzlibrary{decorations.pathreplacing, shapes.geometric, arrows.meta, positioning}

% Niente vedove/orfane: i paragrafi non lasciano righe isolate a fine/inizio pagina
\widowpenalty=10000
\clubpenalty=10000

% Elenchi più compatti
\setlist{itemsep=2pt, parsep=0pt, topsep=4pt, partopsep=0pt}

% --- Riquadri con bordo nero, senza numero (si spezzano tra due pagine) ---
% Uso: \begin{definizione} ... \end{definizione}
%      \begin{teorema}[Nome] ... \end{teorema}  (il nome è facoltativo)
\tcbset{nota/.style={enhanced jigsaw, breakable, colback=white,
  colframe=black, boxrule=0.5pt, sharp corners}}
\NewTColorBox{definizione}{}{nota,
  before upper={\textbf{Definizione.}\ }}
\NewTColorBox{teorema}{o}{nota,
  before upper={\textbf{Teorema\IfValueT{#1}{ (#1)}.}\ }}
\NewTColorBox{proposizione}{o}{nota,
  before upper={\textbf{Proposizione\IfValueT{#1}{ (#1)}.}\ }}
\NewTColorBox{proprieta}{o}{nota,
  before upper={\textbf{Proprietà\IfValueT{#1}{ (#1)}.}\ }}
\NewTColorBox{assioma}{o}{nota, boxrule=1pt,
  before upper={\textbf{Assioma\IfValueT{#1}{ (#1)}.}\ }}

% --- Ambienti semplici (senza riquadro) ---
% Il titolo sta in cima al blocco, anche se il contenuto inizia con un riquadro o
% con codice; e non resta mai da solo in fondo a una pagina.
% Uso: \begin{esempio}[Titolo facoltativo] ... \end{esempio}
\newcommand{\newrunin}[2]{%
  \NewDocumentEnvironment{#1}{o}{%
    \par\Needspace{4\baselineskip}\addvspace{\medskipamount}\noindent
    \textbf{#2}\IfValueT{##1}{ (##1)}\textbf{.}\ \ignorespaces}%
  {\par\addvspace{\medskipamount}}}
\newrunin{esempio}{Esempio}
\newrunin{esercizio}{Esercizio}
\newrunin{osservazione}{Osservazione}

% V e F nelle tavole di verità
\newcommand{\V}{\textsf{V}}
\newcommand{\F}{\textsf{F}}

% --- Codice C ---
% Uso: \begin{codice} ... \end{codice}      codice C numerato
%      \begin{comando} ... \end{comando}    comandi da terminale
%      \begin{risultato} ... \end{risultato} output di un programma
%      \lstinline|printf("ciao");|           codice dentro al testo
\lstdefinestyle{c}{
  language=C,
  basicstyle=\ttfamily\small,
  keywordstyle=\color{blue}\bfseries,
  commentstyle=\color{gray}\itshape,
  stringstyle=\color{red!70!black},
  numbers=left,
  numberstyle=\tiny\color{gray},
  numbersep=8pt,
  columns=flexible,
  keepspaces=true,
  breaklines=true,
  showstringspaces=false,
  tabsize=4,
  morekeywords={include, define},
  % lettere accentate nei commenti
  literate={à}{{\`a}}1 {è}{{\`e}}1 {é}{{\'e}}1 {ì}{{\`i}}1
           {ò}{{\`o}}1 {ù}{{\`u}}1 {È}{{\`E}}1 {À}{{\`A}}1
}
\lstset{style=c}
\newtcblisting{codice}{listing only, nota, left=6mm,
  listing options={style=c}}
\newtcblisting{comando}{listing only, nota,
  listing options={basicstyle=\ttfamily\small, numbers=none, language={},
    columns=flexible, keepspaces=true, breaklines=true}}
\newtcblisting{risultato}{listing only, enhanced jigsaw, breakable,
  colback=black!5, colframe=black, boxrule=0.5pt, sharp corners,
  listing options={basicstyle=\ttfamily\small, numbers=none, language={},
    columns=flexible, keepspaces=true, breaklines=true}}

% --- Diagrammi di flusso (simboli standard) ---
\tikzset{
  terminale/.style = {ellipse, draw, minimum width=2.5cm, minimum height=0.9cm},
  io/.style        = {trapezium, trapezium left angle=70, trapezium right angle=110,
                      draw, minimum width=2.5cm, minimum height=0.9cm},
  processo/.style  = {rectangle, draw, minimum width=2.5cm, minimum height=0.9cm},
  decisione/.style = {diamond, draw, aspect=2, minimum width=2.5cm},
  freccia/.style   = {thick, -{Stealth}}
}

\title{Principio di induzione}
\author{Marco Cerbella}
\date{Lezione 4 -- 29 settembre 2026}

\begin{document}
\maketitle

\section{Il principio di induzione}

\begin{osservazione}
Sia $U \subseteq \mathbb{N}$ tale che
\begin{enumerate}
    \item $0 \in U$;
    \item $k \in U \Rightarrow k + 1 \in U$.
\end{enumerate}
Allora $U = \mathbb{N}$.
\end{osservazione}

Infatti $0 \in U$, quindi $1 \in U$, quindi $2 \in U$, e così via: ``si sale
a gradini'' e si raggiungono tutti i naturali. Se $p(n)$ è un predicato
sui naturali, si applica questa idea all'insieme $U$ dei $n$ per cui
$p(n)$ è vero.

\begin{teorema}[Principio di induzione]
Sia $p(n)$ un predicato. Supponiamo che:
\begin{enumerate}
    \item $p(0)$ sia vero \hfill [\emph{base dell'induzione}];
    \item $\forall\, k \ \big(p(k) \text{ vero} \Rightarrow p(k+1)
          \text{ vero}\big)$ \hfill [\emph{passo induttivo}].
\end{enumerate}
Allora $\forall\, n \in \mathbb{N}$ $p(n)$ è vero.
\end{teorema}

Nel passo induttivo l'ipotesi $p(k)$ si chiama \emph{ipotesi induttiva} e
$p(k+1)$ è la \emph{tesi da dimostrare}.

\begin{esempio}
Dimostrare che per ogni $n$ naturale $2 \mid n(n+1)$.

Sia $p(n)$ il predicato ``$2 \mid n(n+1)$''.

\textbf{1) Caso base.} $p(0)$: $2 \mid 0 \cdot (0 + 1) = 0$. Infatti
$2 \mid 0$ significa $0 = k \cdot 2$ con $k \in \mathbb{Z}$, ed è vero per
$k = 0$: $0 = 0 \cdot 2$.

\textbf{2) Passo induttivo.} Dobbiamo dimostrare $p(k) \Rightarrow
p(k+1)$.
\begin{itemize}
    \item \emph{Ipotesi induttiva} $p(k)$: $2 \mid k(k+1)$.
    \item \emph{Tesi da dimostrare} $p(k+1)$: $2 \mid (k+1)\big((k+1)+1\big)$.
\end{itemize}
Calcoliamo
\[
(k+1)(k+2) = (k+1)k + 2(k+1).
\]
Per l'ipotesi induttiva $2 \mid k(k+1)$, cioè $k(k+1) = 2m$ con
$m \in \mathbb{Z}$. Quindi
\[
(k+1)(k+2) = 2m + 2(k+1) = 2\,(m + k + 1),
\]
cioè $(k+1)(k+2)$ è multiplo di $2$, ovvero $2 \mid (k+1)(k+2)$, che è la
tesi da dimostrare.
\end{esempio}

\section{Induzione a partire da $\bar n$}

Se il predicato è vero solo da un certo naturale in poi, si parte da
quel punto.

\begin{teorema}
Sia $p(n)$ una proposizione e sia $\bar n \in \mathbb{N}$. Supponiamo che:
\begin{enumerate}
    \item $p(\bar n)$ sia vero;
    \item $\forall\, k \geq \bar n \ \big(p(k) \text{ vero} \Rightarrow
          p(k+1) \text{ vero}\big)$.
\end{enumerate}
Allora $\forall\, n \geq \bar n$ $p(n)$ è vera.
\end{teorema}

\begin{esempio}
Il predicato $p(n)$: ``$n - 5 > 0$'' è vero da $6$ in poi: si applica il
teorema con $\bar n = 6$.
\end{esempio}

\begin{esempio}
Dimostrare che per ogni $n \geq 1$
\[
1 + 2 + \dots + n = \frac{n(n+1)}{2}. \tag*{$p(n)$}
\]

\textbf{1) Caso base} $\bar n = 1$. $p(1)$: $1 = \dfrac{1 \cdot (1+1)}{2}$,
vera.

\textbf{2) Passo induttivo.}
\begin{itemize}
    \item \emph{Ipotesi induttiva} $p(k)$:
          $1 + \dots + k = \dfrac{k(k+1)}{2}$.
    \item \emph{Tesi da dimostrare} $p(k+1)$:
          $1 + \dots + k + (k+1) = \dfrac{(k+1)(k+2)}{2}$.
\end{itemize}
Si ha
\begin{align*}
1 + 2 + \dots + k + (k+1) &= (1 + \dots + k) + (k+1)\\
&= \frac{k(k+1)}{2} + (k+1) \qquad \text{(per ipotesi induttiva)}\\
&= \frac{k(k+1) + 2(k+1)}{2} = \frac{(k+1)(k+2)}{2},
\end{align*}
che è la tesi. Per induzione $p(n)$ vale per ogni $n \geq 1$.
\end{esempio}

\end{document}
